% !TEX TS-program = XeLaTeX
% !TEX spellcheck = en-US
\documentclass[aspectratio=169]{beamer}

\usetheme{example}

\title{Lecture 10:\\ Traditional ML vs Deep Learning}
\institute{GRA4160: Predictive modelling with machine learning}
\date{March 11th 2026}
\author{Vegard H\o ghaug Larsen}

\begin{document}

\maketitle

\frame{
\frametitle{Plan for today:}
\begin{itemize}
	\item Traditional ML vs Deep Learning - Overview
	\item Practical applications and model comparisons
	\item Tools used in practice
	\item Emerging trends in machine learning
	% \item Recap of exercise from last week and some more PyTorch
	%\item Visualization resources
  \item Feedback and queestions on mini project
\end{itemize}}

\frame{
\frametitle{Traditional ML vs Deep Learning}
\textbf{Traditional ML Models:}
\begin{itemize}
	\item Algorithms: $k$-NN, decision trees, logistic regression, SVMs
	\item Manual feature selection required
	\item Interpretable, good for structured data, smaller datasets
	\item Example: Decision trees as a flowchart
\end{itemize}
\textbf{Deep Learning:}
\begin{itemize}
	\item Multi-layer neural networks
	\item Automatic feature learning from data
	\item Excellent with large datasets, images, text, audio
	\item High accuracy, high computational cost, low interpretability
	\item Often "black-box" models
\end{itemize}
}

\frame{
\frametitle{Performance vs Dataset Size}
\begin{itemize}
	\item Traditional ML models plateau quickly --- adding more data beyond a threshold yields diminishing returns
	\item Deep learning models scale better with data: performance continues to improve with larger datasets
	\item With \textbf{small data}, traditional ML often outperforms deep learning due to fewer parameters and less overfitting
	\item With \textbf{large data}, deep learning leverages its capacity to learn richer representations
	\item Key insight: the crossover point depends on task complexity and feature quality
\end{itemize}
\bigskip
\centering
\textit{``More data beats a cleverer algorithm, but only if the algorithm can exploit it.''}
}

\frame{
\frametitle{When to Choose Which?}
\begin{itemize}
	\item \textbf{Traditional ML}:
	\begin{itemize}
		\item Limited data
		\item Interpretability critical (e.g., credit decisions)
	\end{itemize}
	\item \textbf{Deep Learning}:
	\begin{itemize}
		\item Large datasets available
		\item Complex, perceptual tasks (vision, NLP)
		\item Computational resources available
	\end{itemize}
	\item Practice: Start simple (Logistic, Random Forest) then move to DL if needed
\end{itemize}}

\frame{
\frametitle{k-NN vs Neural Networks (Recommenders)}
\textbf{k-NN:}
\begin{itemize}
	\item Simple, memory-based, no explicit training
	\item Good for small-scale recommendations
	\item Slow at scale, limited to direct similarity
\end{itemize}
\textbf{Neural Networks:}
\begin{itemize}
	\item Learns embeddings for user-item interactions
	\item Scalable, handles diverse data, highly accurate
	\item Used by Netflix, YouTube
\end{itemize}}

\frame{
\frametitle{Logistic Regression vs DL (Text Classification)}
\textbf{Logistic Regression:}
\begin{itemize}
	\item Baseline, simple linear model, interpretable
	\item Effective with engineered features (e.g. TF-IDF)
	\item Limited nuance in complex texts
\end{itemize}
\textbf{Deep Learning (CNNs, LSTMs, Transformers):}
\begin{itemize}
	\item Surpassed traditional methods on NLP tasks
	\item Learns word/context embeddings
	\item More data-intensive, computationally demanding
\end{itemize}}

\frame{
\frametitle{Decision Trees vs Neural Networks (Medical)}
\textbf{Decision Trees/Ensembles:}
\begin{itemize}
	\item Highly interpretable, suitable for small datasets
	\item Used for simple diagnostic rules
	\item Easy to validate by clinicians
\end{itemize}
\textbf{Deep Learning:}
\begin{itemize}
	\item Human-level accuracy (X-rays, pathology slides)
	\item Harder to interpret, "black-box"
	\item Explainability (XAI) essential in practice
\end{itemize}}

\frame{
\frametitle{Random Forests vs CNNs (Images)}
\textbf{Random Forests:}
\begin{itemize}
	\item Requires extensive feature engineering
	\item Limited scalability and accuracy
\end{itemize}
\textbf{CNNs:}
\begin{itemize}
	\item Revolutionized image recognition
	\item Learn hierarchical features directly from pixels
	\item State-of-the-art accuracy, computationally demanding
\end{itemize}}

\frame{
\frametitle{Ensemble Methods vs DL (Finance)}
\textbf{Ensemble Methods:}
\begin{itemize}
	\item Good with structured/tabular financial data
	\item Robust with smaller datasets
	\item Interpretable, less tuning required
\end{itemize}
\textbf{Deep Learning:}
\begin{itemize}
	\item Powerful with large historical datasets
	\item Captures subtle, complex interactions
	\item Harder to interpret, extensive tuning required
\end{itemize}}

\frame{
\frametitle{Summary: Traditional ML vs Deep Learning}
\centering
\small
\begin{tabular}{lll}
\hline
\textbf{Aspect} & \textbf{Traditional ML} & \textbf{Deep Learning} \\
\hline
Feature engineering & Manual & Automatic \\
Data requirement & Small--medium & Large \\
Interpretability & High & Low \\
Computational cost & Low & High \\
Structured data & Excellent & Good \\
Unstructured data & Limited & Excellent \\
Training time & Fast & Slow \\
Overfitting risk & Lower & Higher \\
\hline
\end{tabular}
\bigskip

Neither approach dominates --- the best choice depends on the \textbf{data}, the \textbf{task}, and the \textbf{constraints}.
}

\frame{
\frametitle{Tools in Practice}
\textbf{Scikit-learn:}
\begin{itemize}
	\item Traditional ML, quick prototyping, easy to use
	\item CPU-based, integrates with pandas, numpy
\end{itemize}
\textbf{TensorFlow/PyTorch:}
\begin{itemize}
	\item Deep learning frameworks, GPU acceleration
	\item TensorFlow: Production deployments
	\item PyTorch: Research, dynamic graphs, HuggingFace integration
\end{itemize}}

\frame{
\frametitle{Emerging Trends in ML}
\begin{itemize}
	\item \textbf{Hybrid Models}: Combining DL features with interpretable ML
	\item \textbf{Transfer/Self-supervised Learning}: Leveraging large, pre-trained models
	\item \textbf{Foundation Models}: GPT, Stable Diffusion, general-purpose models
	\item \textbf{AutoML/No-Code ML}: Automated model building, increased accessibility
\end{itemize}}

\frame{
\frametitle{Transfer Learning in Practice}
\begin{itemize}
	\item \textbf{Idea}: Take a model pre-trained on a large dataset and adapt it to a new, smaller task
	\item Common workflow:
	\begin{enumerate}
		\item Start with a pre-trained model (e.g., ResNet for images, BERT for text)
		\item Freeze early layers that capture general features
		\item Fine-tune later layers on your specific dataset
	\end{enumerate}
	\item \textbf{Why it works}: Early layers learn generic features (edges, textures, word meanings) that transfer across tasks
	\item Dramatically reduces data requirements and training time
	\item Example: Fine-tuning a pre-trained BERT model for sentiment analysis can achieve strong results with just a few hundred labeled examples
\end{itemize}
}

\frame{
\frametitle{Practical Decision Checklist}
When choosing between traditional ML and deep learning, consider:
\begin{enumerate}
	\item \textbf{Data size}: $<$10k samples? Start with traditional ML
	\item \textbf{Data type}: Tabular $\rightarrow$ tree-based methods; images/text/audio $\rightarrow$ deep learning
	\item \textbf{Interpretability}: Required by regulation or stakeholders? Prefer traditional ML or add XAI methods
	\item \textbf{Compute budget}: Limited GPU access? Traditional ML or smaller DL models
	\item \textbf{Iteration speed}: Need quick experiments? Scikit-learn is faster to prototype
	\item \textbf{Baseline first}: Always start with a simple model --- you need a benchmark to justify complexity
\end{enumerate}
\bigskip
\textit{Rule of thumb: if a gradient boosted tree (XGBoost/LightGBM) works well on your tabular data, deep learning will rarely beat it by much.}
}

% \frame{
% \frametitle{Visualization Resources}
% \textbf{Interactive:}
% \begin{itemize}
% 	\item TensorFlow Playground: \url{https://playground.tensorflow.org}
% 	\item Decision trees, SHAP value visualizations
% \end{itemize}
% \textbf{Static:}
% \begin{itemize}
% 	\item Feature Engineering vs Feature Learning (Arivis blog)
% 	\item "Deep Learning" (Goodfellow et al.) visualizations
% 	\item Distill.pub interactive articles
% \end{itemize}}

\end{document}
